\section{A tesztelési stratégia célja}

A WatchDB fejlesztése során a tesztelés célja az volt, hogy a rendszer legfontosabb működési pontjai ellenőrizhetőek legyenek, és a későbbi módosítások ne okozzanak észrevétlen regressziós hibákat. Egy webalkalmazás esetében a felhasználói felület, az API végpontok, a külső szolgáltatásokból érkező adatok és az adatbázis-műveletek egymásra épülnek, ezért egy kisebb módosítás is több rétegre lehet hatással. A tesztek szerepe ebben a folyamatban az, hogy ezekre a kockázatos kapcsolódási pontokra gyors visszajelzést adjanak.

A tesztelési stratégia elsősorban az MVP szempontjából kritikus funkciókra koncentrál. Ide tartozik a filmkeresés, a TMDB-ből érkező adatok feldolgozása, a watchlist kezelése, a filmek megnézettként való jelölése, az értékelés megadása, valamint az autentikációhoz kötött műveletek helyes kezelése. Ezek azok a részek, amelyek hibája közvetlenül rontaná a felhasználói élményt, vagy hibás adatállapotot eredményezne.

A tesztek több szinten jelennek meg a projektben. Az egyszerűbb unit tesztek az izolált függvényeket és service rétegeket ellenőrzik, például az API válaszok normalizálását vagy a formázó függvényeket. A komponens tesztek a felhasználói felület önálló elemeit vizsgálják, például az értékelő komponenst és a megnézettként jelöléshez tartozó dialógust. Az integrációs jellegű tesztek ezzel szemben már több réteg együttműködését ellenőrzik, például azt, hogy egy gombnyomás hatására a megfelelő API hívás történik-e, és a felület a válasz alapján frissül-e.

Fontos szempont volt, hogy a tesztek ne függjenek valódi külső szolgáltatásoktól. A TMDB API, a Supabase kliens, a böngésző \texttt{fetch} hívásai és bizonyos Next.js funkciók a tesztekben mockolt formában jelennek meg. Ez lehetővé teszi, hogy a tesztek gyorsan, determinisztikusan és internetkapcsolat vagy valós adatbázis nélkül fussanak le.

A tesztelés a CI/CD folyamatba is beépül. A GitHub Actions workflow külön lépésekben futtatja a lintelést, a TypeScript típusellenőrzést, a teszteket és a production buildet. Így a Pull Requestek ellenőrzése automatizált módon történhet, és a hibás vagy nem buildelhető módosítások még a fő ágba való bekerülés előtt kiszűrhetők.
